\documentclass[../main.tex]{subfiles}
\ifSubfilesClassLoaded{\setcounter{chapter}{7}}{}
\begin{document}

\chapter{Larik Satu dan Dua Dimensi}

\begin{subcpmk}
  \item Sub-CPMK 4.1: Larik 1D/2D, iterasi elemen dengan aman (batas indeks)
  \item Sub-CPMK 4.3: Pencarian linear dan pengurutan sederhana (bubble sort)
\end{subcpmk}

\SesiRPS{8}{4.1, 4.3}{$\sim$170 menit}{CPMK-4}

% ============================================================
\section{Sarana dan prasarana}
% ============================================================
Sama bab sebelumnya.

% ============================================================
\section{Konsep Array}
% ============================================================

Array (larik) adalah kumpulan elemen bertipe sama yang disimpan dalam \textbf{blok memori berurutan} (\textit{contiguous}). Setiap elemen diakses melalui \textbf{indeks} yang dimulai dari \texttt{0} (bukan \texttt{1}). Array sangat penting dalam pemrograman karena memungkinkan kita menyimpan dan memproses sekumpulan data menggunakan loop, tanpa harus mendeklarasikan variabel terpisah untuk setiap elemen.

Dalam C, ukuran array statis harus ditentukan saat kompilasi (kecuali VLA di C99). C \textbf{tidak memeriksa batas indeks} secara otomatis --- mengakses elemen di luar batas (\textit{out-of-bounds access}) menyebabkan \textit{undefined behavior}, yang merupakan salah satu sumber bug paling berbahaya dalam C. Menjaga agar indeks selalu dalam rentang $[0, n-1]$ adalah tanggung jawab programmer.

\begin{konsep}
\textbf{Representasi memori array \texttt{int arr[5] = \{10,20,30,40,50\}}:}

\begin{modultable}[]{Ilustrasi indeks dan nilai elemen array 1D}
\begin{tabular}{|c|c|c|c|c|}
\hline
\texttt{arr[0]} & \texttt{arr[1]} & \texttt{arr[2]} & \texttt{arr[3]} & \texttt{arr[4]} \\
\hline
10 & 20 & 30 & 40 & 50 \\
\hline
\multicolumn{5}{|c|}{$\underbrace{\hspace{6cm}}_{5 \times \texttt{sizeof(int)}}$ byte berurutan di memori} \\
\hline
\end{tabular}
\end{modultable}
\end{konsep}

% ============================================================
\section{Array Satu Dimensi}
% ============================================================

\begin{lstlisting}[language=C,caption=Deklarasi dan operasi array 1D]
#include <stdio.h>

int main(void) {
    // Deklarasi dan inisialisasi
    int nilai[5] = {85, 92, 78, 96, 88};
    int n = 5;
    
    // Iterasi: tampilkan semua elemen
    printf("Nilai: ");
    for (int i = 0; i < n; i++) {
        printf("%d ", nilai[i]);
    }
    printf("\n");
    
    // Hitung rata-rata
    int total = 0;
    for (int i = 0; i < n; i++) {
        total += nilai[i];
    }
    double rata = (double)total / n;
    printf("Rata-rata: %.2f\n", rata);
    
    // Cari nilai tertinggi
    int maks = nilai[0];
    for (int i = 1; i < n; i++) {
        if (nilai[i] > maks) {
            maks = nilai[i];
        }
    }
    printf("Nilai tertinggi: %d\n", maks);
    
    return 0;
}
\end{lstlisting}

\lstinputlisting[language=C,caption={\texttt{contoh\_code/c/array\_isi\_cetak.c} --- Mengisi dan mencetak array dengan \texttt{while}+\texttt{for} (konversi Pascal)}]{../contoh_code/c/array_isi_cetak.c}

\begin{modultable}[]{Contoh pola inisialisasi array statis}
\begin{tabular}{|l|p{9cm}|}
\hline
\textbf{Inisialisasi} & \textbf{Keterangan} \\
\hline
\texttt{int a[5] = \{1,2,3,4,5\}} & Ukuran eksplisit, semua elemen diisi \\
\hline
\texttt{int a[] = \{1,2,3\}} & Ukuran otomatis dari jumlah elemen (3) \\
\hline
\texttt{int a[5] = \{1,2\}} & Sisa elemen diisi 0 (\texttt{a = \{1,2,0,0,0\}}) \\
\hline
\texttt{int a[5] = \{0\}} & Semua elemen diisi 0 \\
\hline
\texttt{int a[5]} & \textbf{Tidak diinisialisasi} --- berisi garbage! \\
\hline
\end{tabular}
\end{modultable}

% ============================================================
\section{Array sebagai Parameter Fungsi}
% ============================================================

Ketika array dikirim ke fungsi, yang dikirim sebenarnya adalah \textbf{pointer} ke elemen pertama (array ``decay'' menjadi pointer). Artinya, fungsi dapat mengubah isi array asli. Karena informasi ukuran hilang, kita biasanya mengirim parameter tambahan \texttt{int n} untuk jumlah elemen.

\begin{lstlisting}[language=C,caption=Array sebagai parameter fungsi]
#include <stdio.h>

void cetak_array(const int arr[], int n) {
    for (int i = 0; i < n; i++) {
        printf("%d ", arr[i]);
    }
    printf("\n");
}

int cari_max(const int arr[], int n) {
    int maks = arr[0];
    for (int i = 1; i < n; i++) {
        if (arr[i] > maks) maks = arr[i];
    }
    return maks;
}

int main(void) {
    int data[] = {45, 12, 89, 33, 67};
    int n = sizeof(data) / sizeof(data[0]);  // hitung jumlah elemen
    
    cetak_array(data, n);
    printf("Max: %d\n", cari_max(data, n));
    
    return 0;
}
\end{lstlisting}

\begin{catatan}
\texttt{sizeof(data)/sizeof(data[0])} menghitung jumlah elemen array --- tetapi \textbf{hanya bekerja} pada array yang dideklarasikan di scope yang sama, bukan pada parameter fungsi (di sana, array sudah menjadi pointer).
\end{catatan}

\lstinputlisting[language=C,caption={\texttt{contoh\_code/c/array\_prosedur.c} --- Fungsi \texttt{isi} dan \texttt{cetak} untuk array (konversi Pascal)}]{../contoh_code/c/array_prosedur.c}

% ============================================================
\section{Array Dua Dimensi (Matriks)}
% ============================================================

Array 2D disimpan sebagai ``array of arrays'' dalam memori secara \textit{row-major} (baris demi baris). Deklarasi \texttt{int m[BARIS][KOLOM]} mengalokasikan \texttt{BARIS $\times$ KOLOM} elemen secara berurutan.

\begin{lstlisting}[language=C,caption=Penjumlahan dua matriks $2\times 3$]
#include <stdio.h>

#define BARIS 2
#define KOLOM 3

void cetak_matriks(int m[BARIS][KOLOM]) {
    for (int i = 0; i < BARIS; i++) {
        for (int j = 0; j < KOLOM; j++) {
            printf("%4d", m[i][j]);
        }
        printf("\n");
    }
}

int main(void) {
    int A[BARIS][KOLOM] = {{1, 2, 3}, {4, 5, 6}};
    int B[BARIS][KOLOM] = {{7, 8, 9}, {10, 11, 12}};
    int C[BARIS][KOLOM];
    
    // Penjumlahan matriks
    for (int i = 0; i < BARIS; i++) {
        for (int j = 0; j < KOLOM; j++) {
            C[i][j] = A[i][j] + B[i][j];
        }
    }
    
    printf("Matriks A:\n"); cetak_matriks(A);
    printf("\nMatriks B:\n"); cetak_matriks(B);
    printf("\nA + B:\n"); cetak_matriks(C);
    
    return 0;
}
\end{lstlisting}

\lstinputlisting[language=C,caption={\texttt{contoh\_code/c/matriks\_3x3.c} --- Input dan cetak matriks $3\times 3$ (konversi Pascal)}]{../contoh_code/c/matriks_3x3.c}

% ============================================================
\section{Pencarian Linear}
% ============================================================

Pencarian linear (\textit{linear search}) memeriksa setiap elemen satu per satu dari awal hingga akhir. Algoritma ini sederhana dan bekerja pada array apa pun (terurut maupun tidak), namun memiliki kompleksitas \textbf{O(n)} --- lambat untuk array besar.

\begin{lstlisting}[language=C,caption=Pencarian linear]
#include <stdio.h>

int cari_linear(const int arr[], int n, int target) {
    for (int i = 0; i < n; i++) {
        if (arr[i] == target) {
            return i;  // ditemukan di indeks i
        }
    }
    return -1;  // tidak ditemukan
}

int main(void) {
    int data[] = {45, 12, 89, 33, 67, 21};
    int n = sizeof(data) / sizeof(data[0]);
    int target = 33;
    
    int hasil = cari_linear(data, n, target);
    if (hasil != -1) {
        printf("%d ditemukan di indeks %d\n", target, hasil);
    } else {
        printf("%d tidak ditemukan\n", target);
    }
    
    return 0;
}
\end{lstlisting}

% ============================================================
\section{Pengurutan Sederhana: Bubble Sort}
% ============================================================

Bubble sort membandingkan pasangan elemen bersebelahan dan menukarnya jika tidak terurut. Proses diulang hingga seluruh array terurut. Meskipun sederhana untuk dipahami, bubble sort memiliki kompleksitas \textbf{O(n\textsuperscript{2})} dan tidak efisien untuk data besar --- di dunia nyata, gunakan \texttt{qsort()} dari \texttt{<stdlib.h>}.

\begin{lstlisting}[language=C,caption=Bubble sort ascending]
#include <stdio.h>

void bubble_sort(int arr[], int n) {
    for (int i = 0; i < n - 1; i++) {
        int sudah_terurut = 1;  // optimasi: deteksi awal
        for (int j = 0; j < n - 1 - i; j++) {
            if (arr[j] > arr[j + 1]) {
                // Tukar
                int temp = arr[j];
                arr[j] = arr[j + 1];
                arr[j + 1] = temp;
                sudah_terurut = 0;
            }
        }
        if (sudah_terurut) break;  // sudah terurut, berhenti
    }
}

int main(void) {
    int data[] = {64, 34, 25, 12, 22, 11, 90};
    int n = sizeof(data) / sizeof(data[0]);
    
    printf("Sebelum: ");
    for (int i = 0; i < n; i++) printf("%d ", data[i]);
    
    bubble_sort(data, n);
    
    printf("\nSesudah: ");
    for (int i = 0; i < n; i++) printf("%d ", data[i]);
    printf("\n");
    
    return 0;
}
\end{lstlisting}

% ============================================================
\section{Referensi sesi}
% ============================================================
\begin{itemize}[leftmargin=*]
  \item cppreference, array. \url{https://en.cppreference.com/w/c/language/array}
  \item Programiz, \textit{C Arrays}. \url{https://www.programiz.com/c-programming/c-arrays}
  \item GeeksforGeeks, \textit{Bubble Sort}. \url{https://www.geeksforgeeks.org/bubble-sort-algorithm/}
\end{itemize}

% ============================================================
\begin{aktivitas}
  \item Implementasikan penjumlahan dua vektor (array 1D) panjang tetap.
  \item Transpose matriks $2\times 3$ menjadi $3\times 2$.
  \item Implementasikan pencarian linear pada array dan tampilkan hasilnya.
  \item Urutkan array bilangan menggunakan bubble sort, cetak sebelum dan sesudah.
  \item Perhatikan kerangka latihan tabel berikut (konversi dari Pascal \texttt{tabel\_latihan.pas}) --- lengkapi bagian \texttt{fpangkat} dan \texttt{cetak}:
\end{aktivitas}

\lstinputlisting[language=C,caption={\texttt{contoh\_code/c/tabel\_latihan.c} --- Latihan: fungsi tabel pangkat (kerangka konversi Pascal)}]{../contoh_code/c/tabel_latihan.c}

\begin{latihan}
  \item Mengapa indeks array C dimulai dari 0?
  \item Apa yang terjadi jika mengakses \texttt{arr[10]} pada array berukuran 5?
  \item Jelaskan mengapa \texttt{sizeof(arr)/sizeof(arr[0])} tidak bekerja di dalam fungsi.
  \item Berapa perbandingan yang dilakukan bubble sort pada array berukuran $n$?
\end{latihan}

\begin{asesmen}
\textbf{Tugas M6:} manipulasi larik (min/max/rata-rata + pencarian linear + bubble sort) dengan uji kasus batas indeks.
\end{asesmen}

\begin{checklist}
  \item Saya menghindari akses di luar batas larik
  \item Saya dapat mengirim array ke fungsi dengan benar
  \item Saya memahami indeks $0 \ldots n-1$
  \item Saya dapat mengimplementasikan pencarian linear
  \item Saya dapat mengimplementasikan bubble sort
\end{checklist}

\begin{rangkuman}
Bab ini membahas array satu dimensi (deklarasi, inisialisasi, iterasi), array dua dimensi (matriks), array sebagai parameter fungsi, pencarian linear O(n), dan bubble sort O(n\textsuperscript{2}).

\textbf{Poin kunci:} indeks $0 \ldots n-1$; array decay menjadi pointer saat dikirim ke fungsi; C tidak cek batas indeks; hitung ukuran dengan \texttt{sizeof} hanya di scope deklarasi.

\textbf{Kata kunci:} array, indeks, iterasi, matriks, pencarian linear, bubble sort, \texttt{sizeof}
\end{rangkuman}

\end{document}
